% ===== 原文 pp.80–83（PDF p.90–93） =====
% [原文 p.80]
\section{\texorpdfstring{$t=0$、$G=U(1)$}{t=0、G=U(1)} 时的面算子}

在 §5.2 中我们证明了：$t=0$ 时 GL-扭转理论的 3d 约化与 Vafa-Witten 理论的 3d 约化完全相同。因此，这两种理论中面算子的分类完全相同。下面我们考虑 3d 约化理论在阿贝尔规范群情形下的边界条件，先处理规范扇区，再处理物质扇区。

% [Tutor Note: p.80 开头的这一等同性是本节全部结论的基础：正因为 $t=0$ 时 GL-扭转理论与 Vafa-Witten 理论在圆上约化后给出相同的 3d 理论（场含量、BPS 方程与作用量都相同，见 §5.2），两者的面算子分类才相同。因此本节关于面算子的结论可直接读作 Vafa-Witten 理论中面算子的结论，反之亦然。]

\subsection{规范扇区}

与 $t=1$ 的情形一样，我们可以对规范场取 Dirichlet 或 Neumann 条件，然后由 BRST 不变性决定其余部分。边界线算子的范畴通过在带相应边界条件的区间上紧化理论、并分析所得 2d TFT 中的膜来确定。

在 Neumann 情形，有效的 2d TFT 是 A-型 2d 规范理论，与 $t=1$ 时相同。如前所述，其边界条件的范畴等价于 $D^b(\mathrm{Coh}(\C[2]))$。

在 Dirichlet 情形，有效的 2d TFT 是一个带两个玻色场的拓扑 $\sigma$-模型：$A_4$ 与 3d 规范场 $A$ 沿区间的和乐。两者都是周期标量，因此 $\sigma$-模型的靶空间是 $T^2$。BPS 方程约化为全纯瞬子方程
\[
  \dd A_3+\star\dd A_4=0,
\]
这意味着我们面对的是以 $T^2$ 为靶空间的 A-模型。其膜的范畴是 $T^2$ 的 Fukaya-Floer 范畴，它相当不平凡（并且由镜像对称等价于椭圆曲线上凝聚层的上有界导出范畴）。A-模型依赖于 $T^2$ 上的辛形式，而后者可以从作用量的拓扑部分读出。令 $A_1$ 与 $A_2$ 为零、并在长度为 $2\pi$ 的区间上约化，它变成
\[
  -\frac{4\pi^2}{e^2}\int_{M_2}\dd A_3\wedge\dd A_4 .
\]
我们可以把这个表达式看作一个辛 2-形式
\[
  \frac{4\pi^2}{e^2}\,\dd x\wedge\dd y
\]
在 2-环面上的拉回积分，该环面带有周期坐标 $x$、$y$，两者周期均为 1。这个 2-环面的辛面积是 $4\pi^2/e^2$。

我们不知道如何描述这个范畴上由边界线算子融合产生的幺半结构。

% [Tutor Note: p.80 的辛面积 $4\pi^2/e^2$ 可用两处结果交叉核对：(i) 作用量的拓扑部分为 $-\frac{4\pi^2}{e^2}\int_{M_2}\dd A_3\wedge\dd A_4$，而 $A_3,A_4$ 的周期均为 1，故 $T^2$ 上 $\int\dd x\wedge\dd y=1$，辛形式系数即 $4\pi^2/e^2$；(ii) p.82 在电磁对偶下 $e^2\mapsto4\pi^2/e^2$ 且 $\cS\mapsto\cS'=4\pi^2/\cS$，代入 $\cS=4\pi^2/e^2$ 得 $\cS'=e^2$，与 $e^2\mapsto4\pi^2/e^2$ 后的面积一致。]

% [原文 p.81]
\subsection{物质扇区}

与 $t=1$ 时一样，我们可以对标量 $\phi_3$ 与 $\phi_4$ 取 Dirichlet 或 Neumann 条件（BRST 不变性要求它们属于同一类型）。在 Dirichlet 情形，区间上的约化给出修正 A-模型，其唯一的玻色场是二维中的一个实 1-形式 $\phi$。如前所述，它的膜的范畴与平凡 TFT 的情形相同，即等价于 $D^b(\mathrm{Coh}(\bullet))$。与 $t=1$ 的情形不同，这里没有“缠绕扇区”，因为标量 $\phi_3$ 与 $\phi_4$ 不是周期的。所以这一情形中边界线算子的范畴是 $D^b(\mathrm{Coh}(\bullet))$，配以其标准的幺半结构。

% [Tutor Note: p.81 原文 “it category of branes” 是笔误，应为 “its category of branes”（其膜的范畴）。依据：同段上一句主语为 “the modified A-model ... whose only bosonic field ...”，此处需要一个物主限定词；译文按正确形式处理。]

若 $\phi_3$ 与 $\phi_4$ 满足 Neumann 条件，则 1-形式 $\phi$ 在 2d 边界上的限制必须为零。在区间上约化，我们得到唯一玻色场为 $\phi_3$ 与 $\phi_4$ 的 A-模型，即以 $T^*\R$ 为靶空间的 A-模型。其膜的范畴是 $T^*\R$ 的 Fukaya-Floer 范畴。由于这应当被看作某个 3d TFT 中边界线算子的范畴，它应当具有幺半结构。既然与 $t=1$ 的物质扇区相比唯一的差别是 $\phi_4$ 的非紧性，我们预期在应用 [42] 的等价之后，这个幺半结构成为 $\R$ 的可构造导出范畴上的标准幺半结构。

% [Tutor Note: p.81 “唯一差别是 $\phi_4$ 的非紧性”为何导致标准幺半结构：$t=1$ 时 $\phi_4$ 是周期标量，可出现缠绕扇区，从而范畴带 $\C^*$-等变结构；$t=0$ 时 $\phi_3,\phi_4$ 都取值于 $\R$，没有整数缠绕标号，故经 [42] 的等价（$T^*\R$ 的 Fukaya-Floer 范畴 $\simeq$ $\R$ 的可构造导出范畴）后，幺半结构就是层复形张量积给出的标准结构。]

\subsection{把各扇区合在一起}

\subsubsection{DD 条件}

DD 边界条件对应于这样的面算子：1-形式 $\phi$ 具有如下形式的固定奇异性
\[
  \beta\frac{\dd r}{r}-\gamma\,\dd\theta,
\]
而规范场 $A$ 的和乐允许涨落，标量场 $\sigma$ 在插入面上为零。为恰当地定义这样的算子，必须切除插入面的一个管状邻域，并在新产生的边界上施加适当条件。

由于 Dirichlet 情形下物质扇区没有有趣的边界条件，边界线算子的范畴与规范扇区中相同，即以 $T^2$ 为靶空间、辛面积为 $\mathfrak{S}=4\pi^2/e^2$ 的 Fukaya-Floer 范畴。从 4d 的观点看，这就是该面算子上线算子的范畴。

电磁对偶把 DD 条件映到它自身。事实上，它不影响物质扇区；而在规范扇区中，它把周期标量 $A_4$ 映为规范场，并把规范场映为周期标量。由于在 DD 情形中 $A_4$ 满足 Neumann 条件，对偶规范场满足 Dirichlet 条件。反之，规范场的 Dirichlet 条件被对偶映为新周期标量的 Dirichlet 条件。对偶的唯一效果是把 $e^2$ 替换为 $4\pi^2/e^2$。

% [原文 p.82]
因此 $T^2$ 的辛面积也被反转：
\[
  \mathfrak{S}\mapsto\mathfrak{S}'=\frac{4\pi^2}{\mathfrak{S}}.
\]
辛面积如上式相联系的两个环面的 Fukaya-Floer 范畴由通常的 T-对偶而等价。此外，我们预期幺半结构（我们尚未确定它！）在 T-对偶下保持不变。

% [Tutor Note: p.82 的映射 $\cS\mapsto4\pi^2/\cS$ 与 $e^2\mapsto4\pi^2/e^2$ 必须一致：$\cS=4\pi^2/e^2$ 时 $4\pi^2/\cS=e^2$，正是把 $e^2$ 换成 $4\pi^2/e^2$ 后的辛面积。若读者把 $\cS'$ 误写成 $4\pi^2 e^2$ 或 $1/\cS$，则该自洽性检验立即失败。]

\subsubsection{NN 条件}

NN 条件对应于这样的面算子：$A$ 具有如下形式的固定奇异性
\[
  \alpha\,\dd\theta,
\]
而 1-形式 $\phi$ 的奇异性允许涨落。为恰当地定义这样的面算子，必须在插入面的管状邻域的边界上施加适当条件。

在两端都取 NN 边界条件的区间上约化，我们得到一个 2d TFT，它是 A-型 2d 规范理论与以 $T^*\R$ 为靶空间的 A-模型的乘积。其膜的范畴是 $T^*\R$ 的 Fukaya-Floer 范畴的等变版本。上面曾猜测它等价于 $\R$ 的等变可构造导出范畴，配以标准幺半结构（导出张量积）。

电磁对偶把 NN 条件映到它自身，理由与 DD 情形相同。它平凡地作用在线算子的范畴上，因为区间约化后存活下来的玻色场（即 $\sigma$、$\phi_3$ 与 $\phi_4$）都不参与对偶。

\subsubsection{DN 条件}

DN 条件对应于这样的面算子：$A$ 与 $\phi$ 都被允许具有涨落的奇异性，而 $\sigma$ 必须在该面算子处为零。在两端都取 DN 边界条件的区间上约化，我们得到以 $T^2$ 为靶空间的 A-模型与以 $T^*\R$ 为靶空间的 A-模型的乘积。其膜的范畴是 $T^2\times T^*\R$ 的 Fukaya-Floer 范畴。电磁对偶把 DN 条件映到它自身。它在线算子范畴上的作用相当于对 $T^2$ 作 T-对偶（对偶在物质扇区上作用平凡）。幺半结构（我们尚未确定它）必须在 T-对偶下保持不变。

% [Tutor Note: p.82 的两处“我们尚未确定”幺半结构，说明这些 T-对偶下的不变性只是作者的预期而非已证结论（原文用 “we expect”）。读者不应把 NN、DN 两类条件下融合结构的 T-对偶不变性当作定理引用。]

% [原文 p.83]
\subsubsection{ND 条件}

这一情形对应于 Gukov-Witten 面算子：$A$ 的和乐固定，而 1-形式 $\phi$ 具有如下形式的固定奇异性
\[
  \beta\frac{\dd r}{r}-\gamma\,\dd\theta.
\]

在两端都取 ND 边界条件的区间上约化，给出一个 2d TFT，它是 A-型 2d 规范理论与一个修正 A-模型的乘积，后者的唯一玻色场是实 1-形式。由于后一理论中没有有趣的边界条件，这一情形中边界条件的范畴与前一理论中相同。也就是说，它是单点上层的 $U(1)$-等变可构造导出范畴，或等价地 $D^b(\mathrm{Coh}(\C[2]))$ [28]。这因此就是坐在 Gukov-Witten 面算子上的线算子的范畴。幺半结构是标准的（导出张量积）。

特别地，由于平凡面算子是 Gukov-Witten 面算子的特例，我们得出结论：$t=0$ 时 GL-扭转理论中体（bulk）线算子的范畴是 $D^b(\mathrm{Coh}(\C[2]))$。用 4d 的话说，这可以解释为：所有体线算子都可以通过取平凡线算子的若干份拷贝之和、并利用 BRST 不变场 $\sigma$ 及其幂的后代（descendant）对它作变形来构造。这与 [35] 的结果一致：那里论证了在 $t=0$ 时 Wilson 线算子与 't Hooft 线算子都不被允许。

电磁对偶把 ND 条件映到它自身。它平凡地作用在线算子的范畴上，因为场 $\sigma$ 不参与对偶。

% [Tutor Note: p.83 “所有体线算子都由平凡线的拷贝经 $\sigma$ 及其幂的后代形变得到”与 [35] 的结论合起来才完整：$t=0$ 时 Wilson 线与 't Hooft 线都不允许，故体线算子范畴不含这两类对象；剩下的变形参数正是 $\sigma$ 的多项式，对应 $\C[2]$ 的坐标，于是体线算子范畴为 $D^b(\mathrm{Coh}(\C[2]))$。可用“平凡面算子是 Gukov-Witten 面算子的特例”这一观察把该结论从面算子化归到体情形。]
